% ===================================================================
% MT-04c: Minitest Verbos en -ar - MUSTERLÖSUNG
% Spanisch Klasse 7
% ===================================================================

\documentclass[11pt, a4paper]{article}

\usepackage[utf8]{inputenc}
\usepackage[T1]{fontenc}
\usepackage[ngerman]{babel}
\usepackage{helvet}
\renewcommand{\familydefault}{\sfdefault}

\usepackage[margin=2cm]{geometry}
\usepackage{enumitem}
\usepackage{tabularx}
\usepackage{booktabs}

\usepackage{xcolor}
\usepackage{tcolorbox}
\tcbuselibrary{skins}

\usepackage{fancyhdr}
\usepackage{lastpage}

% FARBEN
\definecolor{spanisch}{HTML}{c41e3a}
\definecolor{grau}{HTML}{888888}
\definecolor{hellgrau}{HTML}{f5f5f5}
\definecolor{loesung}{HTML}{c62828}

\newcommand{\fachfarbe}{spanisch}
\newcommand{\thema}{Minitest: Verbos en -ar -- MUSTERLÖSUNG}
\newcommand{\fach}{Spanisch}
\newcommand{\klasse}{7}
\newcommand{\maxpunkte}{28}
\newcommand{\zeit}{15 Minuten}
\newcommand{\datum}{\today}

\pagestyle{fancy}
\fancyhf{}
\renewcommand{\headrulewidth}{0.4pt}
\renewcommand{\footrulewidth}{0.4pt}

\lhead{\fach{} -- Klasse \klasse}
\chead{\textbf{MT-04c -- ML}}
\rhead{\datum}

\lfoot{\zeit}
\cfoot{}
\rfoot{Seite \thepage\ von \pageref{LastPage}}

\newcommand{\punktebox}[1]{\hfill\fbox{\small \_\_\_/#1}}
\newcommand{\luecke}[1]{\rule{#1}{0.4pt}}
\newcommand{\lsg}[1]{\textcolor{loesung}{\textbf{#1}}}

\newtcolorbox{infobox}{
    colback=hellgrau,
    colframe=\fachfarbe,
    boxrule=1pt,
    arc=0pt,
    left=8pt, right=8pt, top=8pt, bottom=8pt
}

\newenvironment{aufgabe}[1]{%
    \vspace{12pt}
    \noindent\textbf{\textcolor{\fachfarbe}{Aufgabe #1}}
    \vspace{6pt}
    \begin{list}{}{\leftmargin=0pt}
    \item[]
}{%
    \end{list}
}

\newcommand{\es}[1]{\textcolor{\fachfarbe}{\textbf{#1}}}
\newcommand{\verbend}[1]{\textcolor{\fachfarbe}{\textbf{#1}}}

\begin{document}

% ===================================================================
% KOPFBEREICH
% ===================================================================

\begin{center}
    {\Large\bfseries\textcolor{loesung}{\thema}}
\end{center}

\vspace{5pt}

\noindent
\begin{tabularx}{\textwidth}{|X|X|X|}
\hline
\textbf{Name:} \lsg{LÖSUNG} &
\textbf{Klasse:} \klasse &
\textbf{Datum:} \luecke{2cm} \\
\hline
\end{tabularx}

\vspace{10pt}

\begin{infobox}
\textbf{Zeit:} \zeit{} | \textbf{Punkte:} \maxpunkte{} BE | \textbf{Hilfsmittel:} keine
\end{infobox}

\vspace{15pt}

% ===================================================================
% AUFGABE 1: Endungen zuordnen
% ===================================================================

\begin{aufgabe}{1 -- Endungen der -ar Verben}
Schreibe die richtige Endung für jede Person. \punktebox{6}
\end{aufgabe}

\noindent
\begin{tabular}{rlrl}
a) yo habl\lsg{-o} & & b) tú estudi\lsg{-as} & \\[0.8em]
c) él/ella escuch\lsg{-a} & & d) nosotros trabaj\lsg{-amos} & \\[0.8em]
e) ellos cant\lsg{-an} & & f) usted bail\lsg{-a} & \\
\end{tabular}

\textcolor{loesung}{\textit{(6 BE: je 1 BE pro richtige Endung)}}

\textcolor{loesung}{\textit{Merke: yo = -o, tú = -as, él/ella/usted = -a, nosotros = -amos, ellos/ellas/ustedes = -an}}

\vspace{15pt}

% ===================================================================
% AUFGABE 2: Konjugation ergänzen
% ===================================================================

\begin{aufgabe}{2 -- Konjugiere die Verben}
Schreibe die konjugierte Form des Verbs in die Lücke. \punktebox{8}
\end{aufgabe}

\noindent
a) (hablar, yo) Yo \lsg{hablo} español.\\[0.8em]
b) (estudiar, tú) Tú \lsg{estudias} mucho.\\[0.8em]
c) (escuchar, María) María \lsg{escucha} música.\\[0.8em]
d) (trabajar, nosotros) Nosotros \lsg{trabajamos} en grupo.\\[0.8em]
e) (cantar, ellos) Ellos \lsg{cantan} en el coro.\\[0.8em]
f) (bailar, yo) Yo \lsg{bailo} salsa.\\[0.8em]
g) (hablar, ellas) Ellas \lsg{hablan} inglés.\\[0.8em]
h) (estudiar, él) Pedro \lsg{estudia} historia.\\

\textcolor{loesung}{\textit{(8 BE: je 1 BE pro korrekte Verbform)}}

\vspace{15pt}

% ===================================================================
% AUFGABE 3: Satzbausteine verbinden
% ===================================================================

\begin{aufgabe}{3 -- Satzbausteine verbinden}
Wähle das passende Verb in der richtigen Form. (A, B, C oder D) \punktebox{8}
\end{aufgabe}

\noindent
\begin{tabular}{|p{7.5cm}|p{5cm}|}
\hline
\textbf{Satz} & \textbf{Optionen} \\
\hline
1. Yo \_\_\_ español en clase. \lsg{A} & A) hablo, B) hablas, C) habla, D) hablamos \\[0.5em]
2. María \_\_\_ matemáticas. \lsg{C} & A) estudio, B) estudias, C) estudia, D) estudiamos \\[0.5em]
3. Nosotros \_\_\_ música en clase. \lsg{D} & A) escucho, B) escuchas, C) escucha, D) escuchamos \\[0.5em]
4. Tú \_\_\_ muy bien. \lsg{B} & A) bailo, B) bailas, C) baila, D) bailamos \\
\hline
\end{tabular}

\textcolor{loesung}{\textit{(8 BE: je 2 BE pro richtige Zuordnung)}}

\vspace{15pt}

% ===================================================================
% AUFGABE 4: Übersetzung
% ===================================================================

\begin{aufgabe}{4 -- Übersetze ins Spanische}
Schreibe die spanische Übersetzung. \punktebox{6}
\end{aufgabe}

\noindent
a) Ich spreche Deutsch. \\[0.5em]
\lsg{(Yo) hablo alemán.} \\[1em]

\noindent
b) Wir arbeiten zusammen. \\[0.5em]
\lsg{(Nosotros) trabajamos juntos.} \\[1em]

\noindent
c) Sie (Plural) singen. \\[0.5em]
\lsg{Ellos cantan. / Ellas cantan.} \\

\textcolor{loesung}{\textit{(6 BE: je 2 BE pro korrekten Satz)}}

% ===================================================================
% PUNKTEÜBERSICHT
% ===================================================================

\vspace{20pt}
\hrule
\vspace{10pt}

\noindent
\begin{tabularx}{\textwidth}{|X|c|c|c|c||c|}
\hline
\textbf{Aufgabe} & 1 & 2 & 3 & 4 & \textbf{Gesamt} \\
\hline
\textbf{max. Punkte} & 6 & 8 & 8 & 6 & \textbf{28} \\
\hline
\textbf{erreicht} & & & & & \\[0.8em]
\hline
\end{tabularx}

\vspace{15pt}

\noindent\textbf{Bewertungsschlüssel (14-NP):}

\begin{center}
\begin{tabular}{|c|c|c|c|c|c|c|c|c|c|c|c|c|c|c|}
\hline
\textbf{NP} & 14 & 13 & 12 & 11 & 10 & 9 & 8 & 7 & 6 & 5 & 4 & 3 & 2 & 1 \\
\hline
\textbf{\%} & 100 & 96 & 91 & 86 & 80 & 73 & 67 & 60 & 53 & 47 & 40 & 33 & 27 & 20 \\
\hline
\textbf{BE} & 28 & 27 & 25 & 24 & 22 & 20 & 19 & 17 & 15 & 13 & 11 & 9 & 8 & 6 \\
\hline
\end{tabular}
\end{center}

\vspace{10pt}

\noindent\textcolor{loesung}{\textbf{Übersicht -ar Konjugation:}}

\begin{center}
\begin{tabular}{|c|c|l|}
\hline
\textbf{Person} & \textbf{Endung} & \textbf{Beispiel (hablar)} \\
\hline
yo & -o & hablo \\
tú & -as & hablas \\
él/ella/usted & -a & habla \\
nosotros/-as & -amos & hablamos \\
ellos/ellas/ustedes & -an & hablan \\
\hline
\end{tabular}
\end{center}

\end{document}
